% Original teaching diagram from the exact example in Chapter 19.
\begin{tikzpicture}[x=1mm,y=1mm,font=\figurefont\fontsize{8}{10}\selectfont,
 v/.style={circle,draw=FigureStroke,fill=white,minimum size=4.5mm,inner sep=0pt},
 e/.style={draw=FigureBlue,line width=.9pt},
 d/.style={draw=FigureGreen,line width=1.1pt},
 sub/.style={font=\FigureNoteFont\fontsize{7.5}{9}\selectfont}]
  \foreach \i/\j/\edgefirst/\edgesecond/\edgethird in {
    0/0/12/23/34,1/0/12/23/41,2/0/12/34/41,3/0/23/34/41,
    0/1/12/34/13,1/1/12/41/13,2/1/23/34/13,3/1/23/41/13} {
    \begin{scope}[shift={({\i*23},{-\j*23})}]
      \coordinate (1) at (0,12);
      \coordinate (2) at (12,12);
      \coordinate (3) at (12,0);
      \coordinate (4) at (0,0);
      \foreach \u/\v/\k in {1/2/12,2/3/23,3/4/34,4/1/41,1/3/13} {
        \ifnum\k=\edgefirst\relax\draw[e] (\u)--(\v);\fi
        \ifnum\k=\edgesecond\relax\draw[e] (\u)--(\v);\fi
        \ifnum\k=\edgethird\relax
          \ifnum\k=13\relax\draw[d] (\u)--(\v);\else\draw[e] (\u)--(\v);\fi
        \fi
      }
      \foreach \k in {1,2,3,4} \node[v] at (\k) {$\k$};
    \end{scope}
  }
  \node[sub] at (40.5,-7) {(a) 不含对角边};
  \node[sub] at (40.5,-30) {(b) 含对角边};
\end{tikzpicture}
